\section{为什么终端会重新成为 AI 入口}

\subsection{终端不是旧界面，而是最小可组合执行层}

如果把 IDE 看作代码编辑中心、浏览器看作信息消费中心，那么终端更像是工程世界的``执行平面''（execution plane）。它有三个很难替代的优势：

\begin{itemize}
    \item \textbf{天然可编排}：命令、管道、环境变量、脚本本身就是可组合积木
    \item \textbf{覆盖面广}：构建、测试、部署、排障、数据处理都能落在命令行上
    \item \textbf{便于代理化}：Agent 可以把动作落成可回放、可审计、可重试的命令序列
\end{itemize}

\begin{knowledgebox}{终端与 Agent 的契合点}
对 Agent 来说，终端并不只是一个黑框 UI，而是一个带有明确输入、输出、权限边界和错误信号的操作系统接口。相比完全图形化的交互，命令行更适合让模型进行计划、执行、观察和修正的循环。
\end{knowledgebox}

\subsection{从 ``autocomplete'' 到 ``delegate'' 的交互跃迁}

早期 AI 编程产品以补全（autocomplete）和问答（chat assist）为主，但 Zach Lloyd 强调，下一阶段的价值来自委托（delegate）：

\begin{enumerate}
    \item 开发者先描述目标，而不是逐行描述实现
    \item 工具生成计划并展示可能采取的动作
    \item 用户在关键节点进行确认，而不是盯着每一步
    \item 系统把结果、错误、日志、diff 再反馈给开发者
\end{enumerate}

\begin{importantbox}{开发者工具的竞争维度在上移}
竞争不再只是``谁补全得更准''，而是``谁能把更高层的工程意图安全地下沉成执行动作''。终端之所以重要，是因为它最接近实际执行，又比 GUI 更容易标准化。
\end{importantbox}

\subsection{本章小结}

终端重新变得关键，并不是因为它复古，而是因为它兼具表达力、可编排性和可审计性。AI 工具若要从建议系统进化成执行系统，终端会是最自然的落点之一。

